% RESEARCH_DRAFT_SHA256: bfe7dc24e2114207de16d40a282c3aa44c9f060cd1892300b0db9c24af268d30
% Layout and source-reference wrappers expanded for the web renderer.
\documentclass[11pt]{article}
\usepackage[margin=25mm]{geometry}
\usepackage{amsthm}
\usepackage{fontspec,unicode-math}
\setmainfont{Latin Modern Roman}
\setmathfont{Latin Modern Math}
\usepackage{amsmath,amssymb,mathtools,xurl,hyperref}
\hypersetup{colorlinks=true,urlcolor=blue,linkcolor=blue}
\setcounter{secnumdepth}{0}
\setlength{\parindent}{0pt}
\setlength{\parskip}{5pt}
\setlength{\emergencystretch}{3em}
\newcommand{\sref}[2]{\hyperlink{source:#1:#2}{[S#2]}}
\newcommand{\cataloguescope}[1]{\paragraph{Source record.}#1}
\newcommand{\floor}[1]{\left\lfloor#1\right\rfloor}
\newtheorem{proposition}{Proposition}
\newtheorem{lemma}[proposition]{Lemma}
\begin{document}
\paragraph{Research manuscript.} OpenAI \texttt{gpt-6-astra}, reasoning \texttt{ultra}; 11 October 2026.
\paragraph{Preserved source body.} The complete catalogue mathematical text and sources follow with layout and source links expanded, including source qualifications and its historical preparation notes. New research begins at the section headed Research attempt.

% UnsolvedMath ID 30005894; source catalogue code OWR-14298366-006
\section{Quantum Harmonic Analysis of Data Structure}\label{30005894}
{\small UnsolvedMath ID 30005894 · Computer Science}
\subsection{Definitions and mathematical statement}

\paragraph{Shared notation.}Standard number systems are $\mathbb N,\mathbb Z,\mathbb Q,\mathbb R,\mathbb C$. Specialized notation and index ranges follow the cited source and the definitions below.


Harmonic analysis decomposes functions or operators into frequency or representation components. In quantum settings operators act on Hilbert spaces and may not commute.
\paragraph{Record scope and qualifications.}
The catalogue title is retained as metadata. Named certificates, contracts, bounds or frameworks in that title are not asserted as results of the original source extract.
Neither the data model nor the mathematical meaning of explained structure is specified.

Source-record limitation: Quantum harmonic analysis of data structure is requested without a mathematical model or target.
\paragraph{Statement recorded in the supplied source.}
Can quantum harmonic analysis explain structure in data? \href{https://doi.org/10.4171/owr/2024/21}{[S2]}
\paragraph{Source status.}
The supplied catalogue labels this source record open; that metadata was checked against the supplied snapshot on 2026-10-11. The label is preserved for provenance and is not an independent verification of present open status or mathematical correctness. The catalogue literature assessment is dated 2026-08-22.
\subsection{Short English statement}
Explore whether quantum harmonic analysis can illuminate structure in data.
\subsection{Sources}
\begin{itemize}
\item[S1] Ulam AI and the UnsolvedMath Contributors, UnsolvedMath: A Curated Collection of Open Mathematics Problems, record 30005894 (OWR-14298366-006). Catalogue curation and metadata: CC BY 4.0. The source extract is retained; definitions, English summaries and qualifications are editorial additions.
\url{https://huggingface.co/datasets/ulamai/UnsolvedMath}.
\item[S2] Applied Harmonic Analysis and Data Science (2024), Oberwolfach Reports, DOI 10.4171/owr/2024/21, pp. 1163--1226
\url{https://doi.org/10.4171/owr/2024/21}.
\item[S3] local structure and dimensionality
\url{https://doi.org/10.1016/j.acha.2024.101692}.
\item[S4] quantum harmonic analysis
\url{https://arxiv.org/abs/2509.19474}.
\item[S5] OWR report
\url{https://doi.org/10.4171/OWR/2024/21}.
\item Source-specific Oberwolfach Reports licences, including CC BY-SA 4.0 where applicable, are retained; see THIRD\_PARTY\_NOTICES.md.
\end{itemize}
\paragraph{Preparation and review.}
The original source and UnsolvedMath contributors retain attribution for the source extract. Codex (GPT-6; exact serving revision not exposed) assisted with definitions, English summaries, scope notes and formatting. No new human mathematical authorship or independently verified solution is asserted. Permanent source ID 30005894 and code OWR-14298366-006 checked against the complete supplied registry and current repository. Distinct upstream records, including repeated questions and broad or damaged extracts, are retained for complete coverage of the source CS category. A repeated record is not a new independent conjecture. Complete mathematical-equivalence and current-literature review has not been performed. Preparation and final review follow contribution prompts C and D, with source gaps disclosed. Structural and rendering checks do not establish a new theorem or open status.
\label{end:30005894}

\section{Research attempt}
\subsection{Literature, source context and scope}
The full original catalogue body and source qualifications are preserved above. The relevant primary workshop passages were read in the cached publisher report. Historical claims and subsequent attributed results are distinguished from the mathematical work below. A restricted example, standard lemma or known theorem is not counted as new progress on the entire original target. Missing assumptions and repeated or rhetorical extracts are not silently turned into different solved problems.

\subsection{What a finite data operator retains and loses}
The source uses data operators and time-frequency localization to study effective dimension. We examine the finite-dimensional analogue $S=\sum_i p_i u_i u_i^*$ for unit vectors $u_i\in\mathbb C^d$, $p_i\ge0$, $\sum_i p_i=1$. This is a positive operator of trace one. Its inverse purity $D(S)=1/\operatorname{tr}(S^2)$ is one possible effective dimension; it is explicitly chosen here, not identified with every entropy-based dimension in the source.
\begin{proposition}
Writing $q=\operatorname{rank}S$, one has
\[
\operatorname{tr}(S^2)=\sum_{i,j}p_ip_j|\langle u_i,u_j\rangle|^2,
\qquad1\le D(S)\le q.
\]
Different point ensembles can nevertheless produce exactly the same $S$ and all functions of $S$.
\end{proposition}
\begin{proof}
Multiplying two rank-one operators and taking their trace gives $|\langle u_i,u_j\rangle|^2$, proving the identity. The nonzero eigenvalues $\lambda_1,\ldots,\lambda_q$ are positive and sum to one. Thus $\sum\lambda_j^2\le1$ and $1=(\sum\lambda_j)^2\le q\sum\lambda_j^2$, proving the bounds.
For nonidentification, in $\mathbb R^2$ give equal mass to $e_1,-e_1,e_2,-e_2$. Its data operator is $I_2/2$. Instead give equal mass to the four vectors $(\pm1,\pm1)/\sqrt2$, with independent choices of signs. Their diagonal outer-product entries are $1/2$ and their averaged off-diagonal entries vanish, so the operator is again $I_2/2$. The supports are disjoint, while both operators have purity $1/2$ and effective dimension two. In addition, the coordinate fourth moment $\mathbb E u_1^4$ is $1/2$ for the first ensemble and $1/4$ for the second, so their distributions differ beyond this operator summary.
\end{proof}
The exact same operator can therefore not identify all distributional geometry or all augmentation directions without extra assumptions or additional information. Unitary changes of representation preserve purity by cyclicity of trace, but an arbitrary feature transformation need not do so. This is consistent with, rather than a rejection of, the source's representation-dependent viewpoint.
\subsection{Verification and remaining gaps}
The two operators, pairwise-overlap purity sums and fourth moments were checked exactly. This finite analogue does not determine the full quantum-harmonic localization operator, prove a nonlinear manifold-recovery theorem, or measure downstream predictive usefulness. The failed inference is from an operator summary to unique recovery of the underlying ensemble. No assertion is made that the source algorithm uses only this summary or claims such uniqueness.

\section{References and source attribution}
Primary workshop report: \url{https://doi.org/10.4171/owr/2024/21}. Relevant passages read: 2024/21 pp.1209--1210, data-operator and localization-operator approach to data structure.

All retrievals were checked on 11 October 2026. Original question attribution remains in the preserved Sources section. External theorems are attributed and not claimed as new discoveries or fully verified merely by citation.


\paragraph{Integration review.} Codex (GPT-6; exact serving revision not exposed) reviewed the written arguments and expanded layout and source-reference wrappers. Repository metadata and the production web renderer were checked using Python 3.13 and JavaScript with MathJax 3.2.2. This model-assisted review is not independent human proof verification; no PDF compilation or formal proof checking is asserted.

\section{Completion Estimate}
\noindent\textbf{COMPLETION: 0\%}
No substantiated new progress toward the entire original target is claimed. The worked results and failed route are delimited above; a known result, illustrative model, or auxiliary obstruction is not counted as original-target completion.

\section{Final status and attribution}
\textbf{Attempt status: unresolved.} This describes this research attempt, not a new ruling on the source catalogue's status or on external literature claims.
The mathematics and this self-audit were generated by OpenAI \texttt{gpt-6-astra}, reasoning \texttt{ultra}, on 11 October 2026. Codex (GPT-6; exact serving revision not exposed) orchestrated the batch, read and reviewed its written mathematical arguments, and contributed formatting/repository review. No human mathematical authorship, independent human proof checking, or formal verification is claimed. The historical preparation paragraphs in the reproduced catalogue body concern that body only.

Contribution prompt A was applied to the complete definition and supplied source record, followed by the adversarial checks of prompt B. Local reproduction files are \texttt{batch025/checks.py}, \texttt{computations.json}, and the per-ID audit JSON. They were not public links at preparation time. Numerical checks are finite and do not prove unrestricted results. No LaTeX engine was found locally; compilation was not performed by this research agent.

\end{document}
